% Replaces §4.4 "Mitigations form a discrete safety--utility Pareto frontier".
% Figure file: results/figures/mitigation_frontier_rows.pdf (analysis/plot_mitigation_frontier.py),
% copied to figures/mitigation_frontier_rows.pdf. The rebuild point in cybersecurity is at three
% seeds for four writers and the canonical seed for Inkling until its two remaining runs finish.
\subsection{Mitigations work when they match the point of entry}
\label{sec:mitigations}
Figure~\ref{fig:mitigation-frontier-domains} places four mitigations on safety--utility axes, one domain per column and one population per row, each mitigation joined by an arrow to its own row's baseline. The two provenance filters, source-authority gating and bounded event sourcing, are evaluated on the paper's five writers at three seeds per domain; the two writer-side changes, rebuilding memory from the history every three blocks and the one-line instruction about who may grant authority, on the five Baseten writers at the same three seeds. The filters move left and far down: on the shared population, gating cuts pooled unauthorized submission from 25.3\% to 7.3\% and event sourcing to 9.0\%, while authorized use falls from 93.3\% to 53.8\% and 64.7\%. Rebuilding moves left with no loss in any domain. The instruction moves left at no measured cost in procurement (unauthorized submission 25.4\% to 5.8\%, paired change $-19.5$ points, 95\% CI $-23.1$ to $-16.0$; authorized use 90.6\% to 98.2\%) and finance (31.7\% to 1.7\%, $-30.0$, 95\% CI $-35.8$ to $-24.2$; authorized use unchanged), at every seed and for every writer, and in the closed loop it all but removes the records minted from the agent's own actions (91 to 4 for the paper's Baseten writers) without recovering authorized use.

In cybersecurity the same instruction moves right and down: unauthorized submission rises from 10.4\% to 21.8\% ($+11.4$, 95\% CI $+6.8$ to $+15.8$) and authorized use falls 10.4 points, at every seed and for four of five writers. The judged failures say where. With or without the instruction, nearly every cybersecurity failure sits at the block carrying the duty officer's signed replacement of the permission list, and in most of them both write attempts are rejected as invalid output (212 of 216 baseline failures, 376 of 403 with the instruction); the instruction makes that rejected update about twice as frequent. Gating does nothing in this domain because every cited source is authoritative. The reading these labels suggest is that a rule about whose word counts does not reach a failure in which the update is never written; the runs that would test it, the same comparison with the write failure removed, are not part of this paper. The domain panels carry the result; pooling the three domains nets the cybersecurity harm against the gains elsewhere, so the frontier is read per domain, jointly from unauthorized submission and authorized use.

\begin{figure}[t]
  \centering
  \includegraphics[width=\linewidth]{figures/mitigation_frontier_rows.pdf}
  \caption{\textbf{Four mitigations on safety--utility axes, by domain and population.} Unauthorized submission against authorized use under typed incremental memory. Top row: the paper's five writers at three seeds per domain, with the two provenance filters, source-authority gating and bounded event sourcing (squares). Bottom row: the five Baseten writers at three seeds per domain, with the two writer-side changes, the one-line instruction and rebuilding from the history every three blocks (diamonds). The circle is each row's own typed-incremental baseline and an arrow runs from it to each mitigation; the rows are different populations and are not compared with each other. Moving left means fewer unauthorized submissions and moving up more authorized use.}
  \label{fig:mitigation-frontier-domains}
\end{figure}
